\section{其他先知书名词}
\subsection{火焰與霹雷、暴風、冰雹}
\subsubsection{以賽亞書-雷轟、地震、大聲、旋風、暴風,火焰 29:5-8}
你仇敵的群眾卻要像細塵，強暴人的群眾也要像飛糠。這事必頃刻之間，忽然臨到。 6 萬軍之耶和華必用\textcolor{red}{雷轟、地震、大聲、旋風、暴風}並吞滅的\textcolor{red}{火焰}，向她討罪。 7 那時，攻擊亞利伊勒列國的群眾，就是一切攻擊亞利伊勒和她的保障並使她困難的，必如夢景，如夜間的異象。 8 又必像飢餓的人，夢中吃飯，醒了仍覺腹空；或像口渴的人，夢中喝水，醒了仍覺發昏，心裡想喝。攻擊錫安山列國的群眾，也必如此。

\subsubsection{以賽亞書-火焰與霹雷、暴風、冰雹 30:27-30}
看哪，耶和華的名從遠方來，怒氣燒起，密煙上騰。他的嘴唇滿有憤恨，他的舌頭像吞滅的火， 28 他的氣如漲溢的河水，直漲到頸項，要用毀滅的篩籮篩淨列國，並且在眾民的口中必有使人錯行的嚼環。 29 你們必唱歌，像守聖節的夜間一樣；並且心中喜樂，像人吹笛，上耶和華的山，到以色列的磐石那裡。 30 耶和華必使人聽他威嚴的聲音，又顯他降罰的膀臂和他怒中的憤恨，並吞滅的\textcolor{red}{火焰與霹雷、暴風、冰雹}。

\subsubsection{啟示錄-火,水變為血,各樣的災殃 11:1-7}
有一根葦子賜給我當做量度的杖，且有話說：「起來，將神的殿和祭壇並在殿中禮拜的人都量一量。 2 只是殿外的院子要留下不用量，因為這是給了外邦人的，他們要踐踏聖城四十二個月。 3 我要使我那兩個見證人穿著毛衣，傳道一千二百六十天。」 4 他們就是那兩棵橄欖樹，兩個燈臺，立在世界之主面前的。 5 若有人想要害他們，就有火從他們口中出來，燒滅仇敵。凡想要害他們的都必這樣被殺。 6 這二人有權柄在他們傳道的日子叫天閉塞不下雨，又有權柄叫水變為血，並且能隨時隨意用各樣的災殃攻擊世界。 7 他們作完見證的時候，那從無底坑裡上來的獸必與他們交戰，並且得勝，把他們殺了。


\subsubsection{啟示錄-地大震動 11:11-13}
過了這三天半，有生氣從神那裡進入他們裡面，他們就站起來，看見他們的人甚是害怕。 12 兩位先知聽見有大聲音從天上來，對他們說：「上到這裡來！」他們就駕著雲上了天，他們的仇敵也看見了。 13 正在那時候，地大震動，城就倒塌了十分之一，因地震而死的有七千人，其餘的都恐懼，歸榮耀給天上的神。

\subsubsection{啟示錄-閃電、聲音、雷轟、地震、大雹 11:15-19}
第七位天使吹號，天上就有大聲音說：「世上的國成了我主和主基督的國，他要做王，直到永永遠遠。」 16 在神面前、坐在自己位上的二十四位長老就面伏於地，敬拜神， 17 說：「昔在、今在的主神，全能者啊！我們感謝你，因你執掌大權做王了。 18 外邦發怒，你的憤怒也臨到了！審判死人的時候也到了！你的僕人眾先知和眾聖徒，凡敬畏你名的人，連大帶小得賞賜的時候也到了！你敗壞那些敗壞世界之人的時候也就到了！」當時，神天上的殿開了，在他殿中現出他的約櫃，隨後有\textcolor{red}{閃電、聲音、雷轟、地震、大雹}。


\subsubsection{啟示錄-閃電、聲音、雷轟、大地震，90斤大雹子 16:12-21}
第六位天使把碗倒在幼發拉底大河上，河水就乾了，要給那從日出之地所來的眾王預備道路。 13 我又看見三個汙穢的靈，好像青蛙，從龍口、獸口並假先知的口中出來。 14 他們本是鬼魔的靈，施行奇事，出去到普天下眾王那裡，叫他們在神全能者的大日聚集爭戰。 15 （「看哪，我來像賊一樣。那警醒、看守衣服，免得赤身而行，叫人見他羞恥的有福了！」） 16 那三個鬼魔便叫眾王聚集在一處，希伯來話叫做哈米吉多頓。

第七位天使把碗倒在空中，就有大聲音從殿中的寶座上出來，說：「成了！」 18 又有\textcolor{red}{閃電、聲音、雷轟、大地震}，自從地上有人以來，沒有這樣大、這樣厲害的地震。 19 那\textcolor{red}{大城裂為三段，列國的城也都倒塌}了。神也想起巴比倫大城來，要把那盛自己烈怒的酒杯遞給他。 20 \textcolor{red}{各海島都逃避了，眾山也不見了}。 21 又有\textcolor{red}{大雹子}從天落在人身上，每一個約重一他連得(一他連得約有九十斤)。為這雹子的災極大，人就褻瀆神。

\subsubsection{啟示錄-火燒淫婦，管轄地上眾王的大城 17:6/14-18}
6我又看見那女人喝醉了聖徒的血和為耶穌作見證之人的血。我看見她，就大大地稀奇。
他們與羔羊爭戰，羔羊必勝過他們，因為羔羊是萬主之主、萬王之王。同著羔羊的，就是蒙召、被選、有忠心的，也必得勝。」 15 天使又對我說：「你所看見那淫婦坐的眾水，就是多民、多人、多國、多方。 16 你所看見的那十角與獸必恨這淫婦，使她冷落赤身，又要吃她的肉，用火將她燒盡。 17 因為神使諸王同心合意，遵行他的旨意，把自己的國給那獸，直等到神的話都應驗了。 18 你所看見的那女人就是管轄地上眾王的大城。

\subsubsection{啟示錄-巴比倫大城傾倒,死亡、悲哀、饑荒,被火燒盡 18:1-24}
此後，我看見另有一位有大權柄的天使從天降下，地就因他的榮耀發光。 2 他大聲喊著說：「巴比倫大城傾倒了！傾倒了！成了鬼魔的住處和各樣汙穢之靈的巢穴(牢獄)，並各樣汙穢可憎之雀鳥的巢穴。 3 因為列國都被她邪淫大怒的酒傾倒了，地上的君王與她行淫，地上的客商因她奢華太過就發了財。」

4 我又聽見從天上有聲音說：「我的民哪，你們要從那城出來，免得與她一同有罪，受她所受的災殃。 5 因她的罪惡滔天，她的不義，神已經想起來了。 6 她怎樣待人，也要怎樣待她，按她所行的加倍地報應她，用她調酒的杯加倍地調給她喝。 7 她怎樣榮耀自己，怎樣奢華，也當叫她照樣痛苦悲哀；因她心裡說：『我坐了皇后的位，並不是寡婦，決不至於悲哀。』 8 所以在一天之內，她的災殃要一齊來到，就是死亡、悲哀、饑荒。她又要被火燒盡了，因為審判她的主神大有能力。 9 地上的君王，素來與她行淫、一同奢華的，看見燒她的煙，就必為她哭泣哀號。 10 因怕她的痛苦，就遠遠地站著說：『哀哉！哀哉！巴比倫大城，堅固的城啊，一時之間你的刑罰就來到了！』 11 地上的客商也都為她哭泣悲哀，因為沒有人再買他們的貨物了。 12 這貨物就是金、銀、寶石、珍珠、細麻布、紫色料、綢子、朱紅色料，各樣香木，各樣象牙的器皿，各樣極寶貴的木頭和銅、鐵、漢白玉的器皿， 13 並肉桂、豆蔻、香料、香膏、乳香、酒、油、細麵、麥子、牛、羊、車、馬和奴僕、人口。

14 「巴比倫哪，你所貪愛的果子離開了你，你一切的珍饈美味和華美的物件也從你中間毀滅，決不能再見了！ 15 販賣這些貨物，藉著她發了財的客商，因怕她的痛苦，就遠遠地站著哭泣悲哀， 16 說：『哀哉，哀哉，這大城啊！素常穿著細麻、紫色、朱紅色的衣服，又用金子、寶石和珍珠為裝飾， 17 一時之間，這麼大的富厚就歸於無有了！』凡船主和坐船往各處去的，並眾水手，連所有靠海為業的，都遠遠地站著， 18 看見燒她的煙就喊著說：『有何城能比這大城呢？』 19 他們又把塵土撒在頭上，哭泣悲哀，喊著說：『哀哉，哀哉，這大城啊！凡有船在海中的都因她的珍寶成了富足，她在一時之間就成了荒場！』 20 天哪，眾聖徒、眾使徒、眾先知啊！你們都要因她歡喜，因為神已經在她身上申了你們的冤。」

21 有一位大力的天使舉起一塊石頭，好像大磨石，扔在海裡說：「巴比倫大城也必這樣猛力地被扔下去，決不能再見了。 22 彈琴、作樂、吹笛、吹號的聲音在你中間決不能再聽見。各行手藝人在你中間決不能再遇見。推磨的聲音在你中間決不能再聽見。 23 燈光在你中間決不能再照耀。新郎和新婦的聲音在你中間決不能再聽見。你的客商原來是地上的尊貴人，萬國也被你的邪術迷惑了。 24 先知和聖徒並地上一切被殺之人的血，都在這城裡看見了。」

\subsection{刀劍、饑荒、瘟疫}
\subsubsection{撒母耳記下 24:13}
於是迦得來見大衛，對他說：「你願意國中有七年的饑荒呢？是在你敵人面前逃跑，被追趕三個月呢？是在你國中有三日的瘟疫呢？現在你要揣摩思想，我好回覆那差我來的。」

\subsubsection{耶 14/15/16/21/24/27/29/32/34/38/42/43/44}
\begin{itemize}

\item 9:16 我要把他們散在列邦中，就是他們和他們列祖素不認識的列邦。我也要使刀劍追殺他們，直到將他們滅盡。」


\item 14:10-16
耶和華對這百姓如此說：「這百姓喜愛妄行 a，不禁止腳步，所以耶和華不悅納他們。現今要記念他們的罪孽，追討他們的罪惡。」 11 耶和華又對我說：「不要為這百姓祈禱求好處。 12 他們禁食的時候，我不聽他們的呼求；他們獻燔祭和素祭，我也不悅納。我卻要用刀劍、饑荒、瘟疫滅絕他們。」13 我就說：「唉，主耶和華啊！那些先知常對他們說：『你們必不看見刀劍，也不遭遇饑荒，耶和華要在這地方賜你們長久的平安。』」 14 耶和華對我說：「那些先知託我的名說假預言，我並沒有打發他們，沒有吩咐他們，也沒有對他們說話。他們向你們預言的，乃是虛假的異象和占卜，並虛無的事以及本心的詭詐。 15 所以耶和華如此說：論到託我名說預言的那些先知，我並沒有打發他們，他們還說『這地不能有刀劍饑荒』，其實那些先知必被刀劍、饑荒滅絕。 16 聽他們說預言的百姓，必因饑荒、刀劍拋在耶路撒冷的街道上，無人葬埋，他們連妻子帶兒女都是如此。我必將他們的惡倒在他們身上。

\item 15:2-3
他們問你說：『我們往哪裡去呢？』你便告訴他們：『耶和華如此說：定為死亡的，必致死亡；定為刀殺的，必交刀殺；定為饑荒的，必遭饑荒；定為擄掠的，必被擄掠。耶和華說：我命定四樣害他們，就是刀劍殺戮，狗類撕裂，空中的飛鳥和地上的野獸吞吃毀滅。

\item 16:4
他們必死得甚苦，無人哀哭，必不得葬埋，必在地上像糞土，必被刀劍和饑荒滅絕。他們的屍首必給空中的飛鳥和地上的野獸做食物。

\item 21:7
以後我要將猶大王西底家和他的臣僕百姓，就是在城內從瘟疫、刀劍、饑荒中剩下的人，都交在巴比倫王尼布甲尼撒的手中和他們仇敵並尋索其命的人手中。巴比倫王必用刀擊殺他們，不顧惜，不可憐，不憐憫。這是耶和華說的。』

\item 24:10
我必使刀劍、饑荒、瘟疫臨到他們，直到他們從我所賜給他們和他們列祖之地滅絕

\item 27:8
『無論哪一邦哪一國，不肯服侍這巴比倫王尼布甲尼撒，也不把頸項放在巴比倫王的軛下，我必用刀劍、饑荒、瘟疫刑罰那邦，直到我藉巴比倫王的手將他們毀滅。這是耶和華說的。

\item 29:16-18 
所以耶和華論到坐大衛寶座的王和住在這城裡的一切百姓，就是未曾與你們一同被擄的弟兄， 17萬軍之耶和華如此說：看哪，我必使刀劍、饑荒、瘟疫臨到他們，使他們像極壞的無花果，壞得不可吃。 18我必用刀劍、饑荒、瘟疫追趕他們，使他們在天下萬國拋來拋去，在我所趕他們到的各國中，令人咒詛、驚駭、嗤笑、羞辱。

\item 32:23-24
他們進入這地得了為業，卻不聽從你的話，也不遵行你的律法。你一切所吩咐他們行的，他們一無所行。因此，你使這一切的災禍臨到他們。 24 看哪，敵人已經來到，築壘要攻取這城，城也因刀劍、饑荒、瘟疫交在攻城的迦勒底人手中。你所說的話都成就了，你也看見了。

\item 32:36-44 現在論到這城，就是你們所說已經因刀劍、饑荒、瘟疫交在巴比倫王手中的，耶和華以色列的神如此說： 37 我在怒氣、憤怒和大惱恨中將以色列人趕到各國，日後我必從那裡將他們招聚出來，領他們回到此地，使他們安然居住。 38 他們要做我的子民，我要做他們的神。 39 我要使他們彼此同心同道，好叫他們永遠敬畏我，使他們和他們後世的子孫得福樂。 40 又要與他們立永遠的約，必隨著他們施恩，並不離開他們，且使他們有敬畏我的心，不離開我。 41 我必歡喜施恩於他們，要盡心、盡意，誠誠實實將他們栽於此地。 42 因為耶和華如此說：我怎樣使這一切大禍臨到這百姓，我也要照樣使我所應許他們的一切福樂都臨到他們。 43 你們說『這地是荒涼無人民、無牲畜，是交付迦勒底人手之地』，日後在這境內必有人置買田地。 44 在便雅憫地、耶路撒冷四圍的各處、猶大的城邑、山地的城邑、高原的城邑並南地的城邑，人必用銀子買田地，在契上畫押，將契封緘，請出見證人，因為我必使被擄的人歸回。這是耶和華說的。

\item 34：17-22 所以耶和華如此說：你們沒有聽從我，各人向弟兄、鄰舍宣告自由，看哪，我向你們宣告一樣自由，就是使你們自由於刀劍、饑荒、瘟疫之下，並且使你們在天下萬國中拋來拋去。這是耶和華說的。 18 猶大的首領，耶路撒冷的首領，太監、祭司和國中的眾民曾將牛犢劈開，分成兩半，從其中經過，在我面前立約。後來又違背我的約，不遵行這約上的話。我必將他們交在仇敵和尋索其命的人手中，他們的屍首必給空中的飛鳥和地上的野獸做食物。 21 並且我必將猶大王西底家和他的首領，交在他們仇敵和尋索其命的人與那暫離你們而去巴比倫王軍隊的手中。 22 耶和華說：我必吩咐他們回到這城，攻打這城，將城攻取，用火焚燒。我也要使猶大的城邑變為荒場，無人居住。

\item 38:2 耶和華如此說：住在這城裡的必遭刀劍、饑荒、瘟疫而死，但出去歸降迦勒底人的必得存活,就是以自己命為掠物的，必得存活。耶和華如此說：這城必要交在巴比倫王軍隊的手中,他必攻取這城。

\item 42:15-22 你們所剩下的猶大人哪，現在要聽耶和華的話！萬軍之耶和華以色列的神如此說：你們若定意要進入埃及，在那裡寄居， 16 你們所懼怕的刀劍在埃及地必追上你們，你們所懼怕的饑荒在埃及要緊緊地跟隨你們，你們必死在那裡。 17 凡定意要進入埃及在那裡寄居的，必遭刀劍、饑荒、瘟疫而死，無一人存留，逃脫我所降於他們的災禍。萬軍之耶和華以色列的神如此說：我怎樣將我的怒氣和憤怒傾在耶路撒冷的居民身上，你們進入埃及的時候，我也必照樣將我的憤怒傾在你們身上，以致你們令人辱罵、驚駭、咒詛、羞辱，你們不得再見這地方。 19 所剩下的猶大人哪，耶和華論到你們說：不要進入埃及去！你們要確實地知道，我今日警教你們了。 20 你們用詭詐自害，因為你們請我到耶和華你們的神那裡，說：『求你為我們禱告耶和華我們神，照耶和華我們的神一切所說的告訴我們，我們就必遵行。』 21 我今日將這話告訴你們，耶和華你們的神為你們的事差遣我到你們那裡說的，你們卻一樣沒有聽從。 22 現在你們要確實地知道：你們在所要去寄居之地必遭刀劍、饑荒、瘟疫而死。

\item 43:11
他要來攻擊埃及地，定為死亡的必致死亡，定為擄掠的必被擄掠，定為刀殺的必交刀殺。

\item 44:12
那定意進入埃及地在那裡寄居的，就是所剩下的猶大人，我必使他們盡都滅絕。必在埃及地仆倒，必因刀劍、饑荒滅絕，從最小的到至大的都必遭刀劍、饑荒而死，以致令人辱罵、驚駭、咒詛、羞辱。

\end{itemize}


耶利米哀歌 1:3
猶大因遭遇苦難，又因多服勞苦，就遷到外邦。她住在列國中，尋不著安息，追逼她的都在狹窄之地將她追上。
耶利米哀歌 4:15
人向他們喊著說：「不潔淨的，躲開！躲開！不要挨近我！」他們逃走漂流的時候，列國中有人說：「他們不可仍在這裡寄居。」



\subsubsection{結 5:12-17/12/14/33}
\begin{itemize}
\item 5:12-17 你的民三分之一必遭\textcolor{red}{瘟疫}而死，在你中間必因\textcolor{red}{饑荒}消滅；三分之一必在你四圍倒在\textcolor{red}{刀}下；我必將三分之一分散四方，並要拔\textcolor{red}{刀}追趕他們。

\item 6:11 主耶和華如此說：你當拍手頓足，說：『哀哉！』以色列家行這一切可憎的惡事，他們必倒在刀劍、饑荒、瘟疫之下。

\item 12:14-16 周圍一切幫助他的和他所有的軍隊，我必分散四方，也要拔刀追趕他們。…我卻要留下他們幾個人得免刀劍、饑荒、瘟疫，使他們在所到的各國中述說他們一切可憎的事，人就知道我是耶和華。

\item 14:13-23 耶和華的話臨到我說： 13 「人子啊，若有一國犯罪干犯我，我也向他伸手折斷他們的杖，就是斷絕他們的糧，使\textcolor{red}{饑荒}臨到那地，將人與牲畜從其中剪除。其中雖有挪亞、但以理、約伯這三人，他們只能因他們的義救自己的性命。這是主耶和華說的。我若使\textcolor{red}{惡獸}經過糟踐那地，使地荒涼，以致因這些獸，人都不得經過， 16 雖有這三人在其中——主耶和華說：我指著我的永生起誓，他們連兒帶女都不能得救，只能自己得救，那地仍然荒涼。或者我使\textcolor{red}{刀劍}臨到那地，說『刀劍哪，要經過那地』，以致我將人與牲畜從其中剪除， 18 雖有這三人在其中——主耶和華說：我指著我的永生起誓，他們連兒帶女都不能得救，只能自己得救。或者我叫\textcolor{red}{瘟疫}流行那地，使我滅命(帶血)的憤怒傾在其上，好將人與牲畜從其中剪除， 20 雖有挪亞、但以理、約伯在其中——主耶和華說：我指著我的永生起誓，他們連兒帶女都不能救，只能因他們的義救自己的性命。主耶和華如此說：我將這四樣大災，就是\textcolor{red}{刀劍、饑荒、惡獸、瘟疫}降在耶路撒冷，將人與牲畜從其中剪除，豈不更重嗎？然而其中必有剩下的人，他們連兒帶女必帶到你們這裡來，你們看見他們所行所為的，要因我降給耶路撒冷的一切災禍，便得了安慰。 23 你們看見他們所行所為的得了安慰，就知道我在耶路撒冷中所行的並非無故。這是主耶和華說的。


\item 33:1-6 耶和華的話臨到我說： 2 「人子啊，你要告訴本國的子民說：我使刀劍臨到哪一國，那一國的民從他們中間選立一人為守望的， 3 他見刀劍臨到那地，若吹角警戒眾民， 4 凡聽見角聲不受警戒的，刀劍若來除滅了他，他的罪(血)就必歸到自己的頭上。 5 他聽見角聲，不受警戒，他的罪必歸到自己的身上；他若受警戒，便是救了自己的性命。 6 倘若守望的人見刀劍臨到，不吹角，以致民不受警戒，刀劍來殺了他們中間的一個人，他雖然死在罪孽之中，我卻要向守望的人討他喪命的罪 b。

\item 33:27-29 你要對他們這樣說：『主耶和華如此說：我指著我的永生起誓，在荒場中的必倒在\textcolor{red}{刀}下，在田野間的必交給\textcolor{red}{野獸}吞吃，在保障和洞裡的必遭\textcolor{red}{瘟疫}而死。我必使這地荒涼，令人驚駭，她因勢力而有的驕傲也必止息。以色列的山都必荒涼，無人經過。 29 我因他們所行一切可憎的事，使地荒涼，令人驚駭，那時他們就知道我是耶和華。』


\end{itemize}


\subsubsection{何西阿書 9:6}
看哪，他們逃避災難，埃及人必收殮他們的屍首，摩弗人必葬埋他們的骸骨。他們用銀子做的美物上必長蒺藜，他們的帳篷中必生荊棘

\subsubsection{阿摩司書 4:6-12}
「我使你們在一切城中牙齒乾淨，在你們各處糧食缺乏，你們仍不歸向我。」這是耶和華說的。
阿摩司書 9:4
雖被仇敵擄去，我必命刀劍殺戮他們。我必向他們定住眼目，降禍不降福。」

撒迦利亞書 7:14
我必以旋風吹散他們到素不認識的萬國中。這樣，他們的地就荒涼，甚至無人來往經過，因為他們使美好之地荒涼了。」

\subsubsection{路加福音 21:24}
他們要倒在刀下，又被擄到各國去。耶路撒冷要被外邦人踐踏，直到外邦人的日期滿了。


\subsubsection{啟示錄 6:4-8}
就另有一匹馬出來，是紅的，有權柄給了那騎馬的，可以從地上奪去太平，使人彼此相殺，又有一把大刀賜給他。
我就觀看，見有一匹灰色馬。騎在馬上的名字叫做「死」，陰府也隨著他。有權柄賜給他們，可以用刀劍、饑荒、瘟疫、野獸殺害地上四分之一的人。


\subsection{屍首必給空中的飛鳥和地上的野獸做食物}
\subsubsection{耶利米書 34：17-22}
所以耶和華如此說：你們沒有聽從我，各人向弟兄、鄰舍宣告自由，看哪，我向你們宣告一樣自由，就是使你們自由於刀劍、饑荒、瘟疫之下，並且使你們在天下萬國中拋來拋去。這是耶和華說的。 18 猶大的首領，耶路撒冷的首領，太監、祭司和國中的眾民曾將牛犢劈開，分成兩半，從其中經過，在我面前立約。後來又違背我的約，不遵行這約上的話。 19  20 我必將他們交在仇敵和尋索其命的人手中，他們的屍首必給空中的飛鳥和地上的野獸做食物。 21 並且我必將猶大王西底家和他的首領，交在他們仇敵和尋索其命的人與那暫離你們而去巴比倫王軍隊的手中。 22 耶和華說：我必吩咐他們回到這城，攻打這城，將城攻取，用火焚燒。我也要使猶大的城邑變為荒場，無人居住。

\subsection{眾星群宿不發光，日頭變黑月亮也不放光,人少}
以賽亞書 13

應當在淨光的山豎立大旗，向群眾揚聲招手，使他們進入貴胄的門！ 3 我吩咐我所挑出來的人，我招呼我的勇士，就是那矜誇高傲之輩，為要成就我怒中所定的。」 4 山間有多人的聲音，好像是大國人民。有許多國的民聚集鬨嚷的聲音，這是萬軍之耶和華點齊軍隊，預備打仗。 5 他們從遠方來，從天邊來，就是耶和華並他惱恨的兵器，要毀滅這全地。你們要哀號，因為耶和華的日子臨近了！這日來到，好像毀滅從全能者來到。 7 所以人手都必軟弱，人心都必消化。 8 他們必驚惶悲痛，愁苦必將他們抓住，他們疼痛好像產難的婦人一樣。彼此驚奇相看，臉如火焰。 9 耶和華的日子臨到，必有殘忍、憤恨、烈怒，使這地荒涼，從其中除滅罪人。 10 天上的眾星群宿都不發光，日頭一出就變黑暗，月亮也不放光。 11 「我必因邪惡刑罰世界，因罪孽刑罰惡人，使驕傲人的狂妄止息，制伏強暴人的狂傲。 12 我必使人比精金還少，使人比俄斐純金更少。 13 我萬軍之耶和華在憤恨中發烈怒的日子，必使天震動，使地搖撼離其本位。 14 人必像被追趕的鹿，像無人收聚的羊，各歸回本族，各逃到本土。 15 凡被仇敵追上的必被刺死，凡被捉住的必被刀殺。 16 他們的嬰孩必在他們眼前摔碎，他們的房屋必被搶奪，他們的妻子必被玷汙。17 「我必激動瑪代人來攻擊他們，瑪代人不注重銀子，也不喜愛金子。 18 他們必用弓擊碎少年人，不憐憫婦人所生的，眼也不顧惜孩子。 19 巴比倫素來為列國的榮耀，為迦勒底人所矜誇的華美，必像神所傾覆的所多瑪、蛾摩拉一樣。 20 其內必永無人煙，世世代代無人居住。阿拉伯人也不在那裡支搭帳篷，牧羊的人也不使羊群臥在那裡。 21 只有曠野的走獸臥在那裡，咆哮的獸滿了房屋，鴕鳥住在那裡，野山羊在那裡跳舞。 22 豺狼必在她宮中呼號，野狗必在她華美殿內吼叫。巴比倫受罰的時候臨近，她的日子必不長久。」

以賽亞書 24
看哪，耶和華使地空虛，變為荒涼；又翻轉大地，將居民分散。 2 那時百姓怎樣，祭司也怎樣；僕人怎樣，主人也怎樣；婢女怎樣，主母也怎樣；買物的怎樣，賣物的也怎樣；放債的怎樣，借債的也怎樣；取利的怎樣，出利的也怎樣。 3 地必全然空虛，盡都荒涼，因為這話是耶和華說的。 4 地上悲哀衰殘，世界敗落衰殘，地上居高位的人也敗落了。 5 地被其上的居民汙穢，因為他們犯了律法，廢了律例，背了永約。 6 所以地被咒詛吞滅，住在其上的顯為有罪；\textcolor{red}{地上的居民被火焚燒，剩下的人稀少}。 7 新酒悲哀，葡萄樹衰殘，心中歡樂的俱都嘆息。 8 擊鼓之樂止息，宴樂人的聲音完畢，彈琴之樂也止息了。 9 人必不得飲酒唱歌，喝濃酒的必以為苦。 10 荒涼的城拆毀了，各家關門閉戶，使人都不得進去。 11 在街上因酒有悲嘆的聲音，一切喜樂變為昏暗，地上的歡樂歸於無有。 12 城中只有荒涼，城門拆毀淨盡。 13 在地上的萬民中，必像打過的橄欖樹，又像已摘的葡萄\textcolor{red}{所剩無幾}。

14 這些人要高聲歡呼，他們為耶和華的威嚴，從海那裡揚起聲來。 15 因此你們要在東方榮耀耶和華，在眾海島榮耀耶和華以色列神的名。 16 我們聽見從地極有人歌唱說：「榮耀歸於義人！」

我卻說：「我消滅了！我消滅了！我有禍了！詭詐的行詭詐，詭詐的大行詭詐。」 17 地上的居民哪，恐懼、陷坑、網羅都臨近你。 18 躲避恐懼聲音的必墜入陷坑，從陷坑上來的必被網羅纏住，因為天上的窗戶都開了，\textcolor{red}{地的根基也震動了。 19 地全然破壞，盡都崩裂，大大地震動了。 20 地要東倒西歪，好像醉酒的人；又搖來搖去，好像吊床}。罪過在其上沉重，必然塌陷，不能復起。


21 到那日，耶和華在高處必懲罰高處的眾軍，在地上必懲罰地上的列王。 22 他們必被聚集，像囚犯被聚在牢獄中，並要囚在監牢裡，多日之後便被討罪 a。 23 那時，\textcolor{red}{月亮要蒙羞，日頭要慚愧}，因為萬軍之耶和華必在錫安山，在耶路撒冷做王，在敬畏他的長老面前必有榮耀。

\subsection{陀斐特}
Tophet(地名)

字義：    焚燒之處 Place of burning，唾液 Spittle，鼓、喧聲 Drum, noise，壇 Altar，輕蔑 on tempt

彙編：    ‧王下23:10；耶7:31；賽30:33是欣嫩子谷中之高阜，以賽亞及耶利米時居民曾使子女經火奉摩洛，約西亞王禁之而污穢其地。另同名相似地乃為亞述王而設。

―― 陳瑞庭

 

陀斐特(Topheth)耶路撒冷城外、欣嫩子谷內的一個地方，以色列人曾在那裏以活人為祭獻給摩洛，褻瀆耶和華（參代下二十八3，三十三6）。

約西亞在其宗教改革中，污穢了陀斐特，並拆毀其中的祭壇（王下二十三10）。約西亞的改革似乎只起了暫時的作用，因後來耶利米對這惡行加以譴責（耶七31、32）。耶利米預言這地方要易名為「殺戮谷」，因為巴比倫人會在圍攻耶路撒冷期間，於這裏大敗猶大國（耶七32）。耶利米講述窯匠之瓦瓶的比喻時，再次重複這個預言。他強調耶路撒冷將要全然被毀，像陀斐特一樣（耶十九12）。那時，陀斐特顯然已成了一個垃圾場，打碎的瓦器會被扔在那裏，無葬身之地的屍首也會搬到陀斐特來（耶十九11）。雖然新約並沒有提及陀斐特，但「欣嫩子谷」可能就是指這個地方，意指一個遭毀壞的地方，新約中常譯作「地獄」（太五22、29、30，十28，十八9；可九43-47；路十二5）。

在加略人猶大賣主的事件中，陀斐特也可能隱晦地被提及。祭司用了猶大賣主得到的銀錢，購買窯戶的一塊田，用以埋葬外鄉人（太二十七9、10；參耶十九1-13，三十二6-9）。一個值得注意的地方是，耶利米書中稱陀斐特為「殺戮谷」，窯戶的田則稱為「血田」（太二十七10；徒一19）。然而，血田確實的位置卻不得而知。―― 證主聖經百科全書

\subsubsection{火與許多木柴早已在陀斐特為王預備 賽30:31-33}
\textcolor{blue}{亞述人必因耶和華的聲音驚惶，耶和華必用杖擊打他。 32 耶和華必將命定的杖加在他身上，每打一下，人必擊鼓彈琴。打仗的時候，耶和華必掄起手來與他交戰。} 33 原來陀斐特又深又寬，早已為王預備好了，其中堆的是火與許多木柴，耶和華的氣如一股硫磺火，使它著起來。

\subsection{二十四位長老 }
\begin{itemize}
\item 啟示錄 4:4 \\
此後，我觀看，見天上有門開了。我初次聽見好像吹號的聲音對我說：「你上到這裡來，我要將以後必成的事指示你。」 2 我立刻被聖靈感動，見有一個寶座安置在天上，又有一位坐在寶座上。 3 看那坐著的，好像碧玉和紅寶石，又有虹圍著寶座，好像綠寶石。寶座的周圍又有二十四個座位，其上坐著\textcolor{red}{二十四位長老}，身穿白衣，頭上戴著金冠冕。
\item 啟示錄 4:9-11\\
 每逢四活物將榮耀、尊貴、感謝歸給那坐在寶座上、活到永永遠遠者的時候， 10 那\textcolor{red}{二十四位長老}就俯伏在坐寶座的面前，敬拜那活到永永遠遠的，又把他們的冠冕放在寶座前，說： 11 「我們的主，我們的神，你是配得榮耀、尊貴、權柄的，因為你創造了萬物，並且萬物是因你的旨意被創造而有的。
\item 啟示錄 5:1-14 我看見坐寶座的右手中有書卷，裡外都寫著字，用七印封嚴了。 2 我又看見一位大力的天使大聲宣傳說：「有誰配展開那書卷，揭開那七印呢？」 3 在天上、地上、地底下，沒有能展開、能觀看那書卷的。 4 因為沒有配展開、配觀看那書卷的，我就大哭。 5 \textcolor{red}{長老中有一位}對我說：「不要哭。看哪，猶大支派中的獅子，大衛的根，他已得勝，能以展開那書卷、揭開那七印！」我又看見寶座與四活物並\textcolor{red}{長老之中}有羔羊站立，像是被殺過的，有七角七眼，就是神的七靈，奉差遣往普天下去的。 7 這羔羊前來，從坐寶座的右手裡拿了書卷。 8 他既拿了書卷，四活物和\textcolor{red}{二十四位長老}就俯伏在羔羊面前，各拿著琴和盛滿了香的金爐。這香就是眾聖徒的祈禱。 9 他們唱新歌，說：「你配拿書卷，配揭開七印，因為你曾被殺，用自己的血從各族、各方、各民、各國中買了人來，叫他們歸於神， 10 又叫他們成為國民做祭司，歸於神，在地上執掌王權。」我又看見且聽見寶座與活物並\textcolor{red}{長老}的周圍有許多天使的聲音，他們的數目有千千萬萬， 12 大聲說：「曾被殺的羔羊是配得權柄、豐富、智慧、能力、尊貴、榮耀、頌讚的！」 13 我又聽見在天上、地上、地底下、滄海裡和天地間一切所有被造之物都說：「但願頌讚、尊貴、榮耀、權勢都歸給坐寶座的和羔羊，直到永永遠遠！」 14 四活物就說：「阿們！」\textcolor{red}{眾長老}也俯伏敬拜。
\item 啟示錄 7:9-17 \\
此後我觀看，見有許多的人，沒有人能數過來，是從各國、各族、各民、各方來的，站在寶座和羔羊面前，身穿白衣，手拿棕樹枝， 10 大聲喊著說：「願救恩歸於坐在寶座上我們的神，也歸於羔羊！」 11 眾天使都站在寶座和\textcolor{red}{眾長老}並四活物的周圍，在寶座前面伏於地，敬拜神， 12 說：「阿們！頌讚、榮耀、智慧、感謝、尊貴、權柄、大力都歸於我們的神，直到永永遠遠！阿們。」\textcolor{red}{长老}中有一位问我：「这些身穿白袍的是谁？他们是从哪里来的呢？」 14我对他说：「我主啊，你是知道的。」他告诉我：「这些人是从大患难中出来的。他们用羊羔的血，把自己的衣袍洗洁白了。 15因此，他们可以在　神的宝座前，并且日夜在他的圣所里事奉他。那坐在宝座上的必展开帐幕覆庇他们。 16他们必不再饿，也不再渴；太阳和一切炎热必不伤害他们。 17因为在宝座中间的羊羔必牧养他们，领他们到生命水的泉源那里。　神也必抹去他们的一切眼泪。」
\item 啟示錄 11:15-18\\
第七位天使吹號，天上就有大聲音說：「世上的國成了我主和主基督的國，他要做王，直到永永遠遠。」 16 在神面前、坐在自己位上的\textcolor{red}{二十四位長老}就面伏於地，敬拜神， 17 說：「昔在、今在的主神，全能者啊！我們感謝你，因你執掌大權做王了。 18 外邦發怒，你的憤怒也臨到了！審判死人的時候也到了！你的僕人眾先知和眾聖徒，凡敬畏你名的人，連大帶小得賞賜的時候也到了！你敗壞那些敗壞世界之人的時候也就到了！」
\item 啟示錄 14:1-5\\
我又觀看，見羔羊站在錫安山，同他又有十四萬四千人，都有他的名和他父的名寫在額上。 2 我聽見從天上有聲音，像眾水的聲音和大雷的聲音，並且我所聽見的好像彈琴的所彈的琴聲。 3 他們在寶座前，並在四活物和\textcolor{red}{眾長老}前唱歌，彷彿是新歌，除了從地上買來的那十四萬四千人以外，沒有人能學這歌。 4 這些人未曾沾染婦女，他們原是童身。羔羊無論往哪裡去，他們都跟隨他。他們是從人間買來的，做初熟的果子歸於神和羔羊。 5 在他們口中察不出謊言來，他們是沒有瑕疵的。
\item 啟示錄 19:1-6\\
此後，我聽見好像群眾在天上大聲說：「哈利路亞 a！救恩、榮耀、權能都屬乎我們的神！ 2 他的判斷是真實、公義的，因他判斷了那用淫行敗壞世界的大淫婦，並且向淫婦討流僕人血的罪，給他們申冤。」 3 又說：「哈利路亞！燒淫婦的煙往上冒，直到永永遠遠！」 4 那\textcolor{red}{二十四位長老}與四活物就俯伏敬拜坐寶座的神，說：「阿們！哈利路亞！」 5 有聲音從寶座出來，說：「神的眾僕人哪，凡敬畏他的，無論大小，都要讚美我們的神！」 6 我聽見好像群眾的聲音，眾水的聲音，大雷的聲音，說：「哈利路亞！因為主我們的神，全能者做王了。
\item 啟示錄 20:4我又看見幾個寶座，也有坐在上面的，並有審判的權柄賜給他們。我又看見那些因為給耶穌作見證並為神之道被斬者的靈魂，和那沒有拜過獸與獸像，也沒有在額上和手上受過牠印記之人的靈魂；他們都復活了，與基督一同做王一千年。 
\end{itemize}

\subsection{番石榴树  הַהֲדַסִּ֖ים }
\subsubsection{尼希米記 8:13-15}
次日，眾民的族長、祭司和利未人都聚集到文士以斯拉那裡，要留心聽律法上的話。 14 他們見律法上寫著，耶和華藉摩西吩咐以色列人要在七月節住棚， 15 並要在各城和耶路撒冷宣傳報告說：「你們當上山，將橄欖樹、野橄欖樹、\underline{番石榴樹}、棕樹和各樣茂密樹的枝子取來，照著所寫的搭棚（And that they should publish and proclaim in all their cities, and in Jerusalem, saying, Go forth unto the mount, and fetch olive branches, and pine branches, and myrtle branches, and palm branches, and branches of thick trees, to make booths, as it is written）。」

\subsubsection{以賽亞書 41:17-20}
17 「困苦窮乏人尋求水卻沒有，他們因口渴舌頭乾燥，我耶和華必應允他們，我以色列的神必不離棄他們。 18 我要在淨光的高處開江河，在谷中開泉源；我要使沙漠變為水池，使乾地變為湧泉。 19 我要在曠野種上香柏樹、皂莢樹、\underline{番石榴樹}和野橄欖樹，我在沙漠要把松樹、杉樹並黃楊樹一同栽植（I will plant in the wilderness the cedar, the shittah tree, and the myrtle, and the oil tree; I will set in the desert the fir tree, and the pine, and the box tree together:）， 20 好叫人看見、知道、思想、明白這是耶和華的手所做的，是以色列的聖者所造的

\subsubsection{以賽亞書 55:8-13}
耶和華說：「我的意念非同你們的意念，我的道路非同你們的道路。 9 天怎樣高過地，照樣，我的道路高過你們的道路，我的意念高過你們的意念。 10 雨雪從天而降並不返回，卻滋潤地土，使地上發芽結實，使撒種的有種，使要吃的有糧。 11 我口所出的話也必如此，決不徒然返回，卻要成就我所喜悅的，在我發它去成就 b的事上必然亨通。 12 你們必歡歡喜喜而出來，平平安安蒙引導，大山小山必在你們面前發聲歌唱，田野的樹木也都拍掌。 13 松樹長出代替荊棘，\underline{番石榴}長出代替蒺藜，這要為耶和華留名，作為永遠的證據，不能剪除（Instead of the thorn shall come up the fir tree, and instead of the brier shall come up the myrtle tree: and it shall be to the LORD for a name, for an everlasting sign that shall not be cut off）」

\subsubsection{撒迦利亞書 1:8-11}
那站在\underline{番石榴樹}中間的人說：「這是奉耶和華差遣在遍地走來走去的。」我對與我說話的天使說：「主啊，這是什麼意思？」他說：「我要指示你這是什麼意思。」 10那站在番石榴樹中間的人說：「這是奉耶和華差遣在遍地走來走去的。」 11那些騎馬的對站在番石榴樹中間耶和華的使者說：「我們已在遍地走來走去，見全地都安息平靜。」
\begin{figure}[h!]
  \centering
  \begin{subfigure}[b]{0.2\linewidth}
    \includegraphics[width=\linewidth]{ZongJie/figs/XianZhi_CreepMyrtle}
     \caption{番石榴樹}
  \end{subfigure}
  \begin{subfigure}[b]{0.2\linewidth}
    \includegraphics[width=\linewidth,height=3.9cm]{ZongJie/figs/XianZhi_CreepMyrtle2}
    \caption{番石榴樹}
  \end{subfigure}
  \begin{subfigure}[b]{0.2\linewidth}
    \includegraphics[width=\linewidth,height=3.85cm]{ZongJie/figs/XianZhi_CreepMyrtle3}
    \caption{番石榴樹}
  \end{subfigure}
\end{figure}



\subsection{窪地  בַּמְּצֻלָ֑ה  ravine,in the bottom}




\subsection{神的榮耀}
榮耀 Glory，希伯來字 Kabowd ＃3519，來自字根 kabad 本意為「重」的意思，引申為價值，尊貴，榮美，豐富，榮譽；以西結在書中，用虹的光彩來描寫神的榮耀。神的榮耀就是祂的本體，彰顯祂的聖潔，榮美，能力，與莊嚴，是神本相的表現。所以詩人說「諸天述說神的榮耀」（詩 19：1），就是藉著神所創造的宇宙萬物，神啟示給我們看到祂的本體。
榮耀的希臘字是 doxa ＃1391，原意是「信念」belief 與「看法」opinion 的意思，由七十士譯本將希伯來文的 Kabowd 譯成希臘文的 doxa，引申為讚美，尊崇，榮耀 ；神的榮耀就是神所有屬性的總和。
%\subsubsection{以西結書}


%\subsubsection{出埃及記}

%摩西在出埃及記中求神彰顯榮耀給他看，神回答「，又說「以西結書從神對選民長久的忍耐（不輕易發怒），回轉的勸誡（有憐憫有恩典），猶大與外邦所受的刑罰（萬不以有罪的為無罪），到餘民復興的預言，聖殿的重建（為千萬人存留慈愛，赦免罪孽），無不彰顯神的榮耀（屬性）。

%\subsubsection{羅 3：23}
%保羅說「世人都犯了罪，虧缺了神的榮耀」（羅 3：23），也就是說我們因著罪的緣故，與神的性情，神的本質分隔。然而神的心意是要人回到神的榮耀中，
%\subsubsection{猶 1：24}
%「叫你們無瑕無疵，歡歡喜喜站在他榮耀之前…」（猶 1：24），
%\subsubsection{啟 21：23}
%「那城內又不用日月光照，因有神的榮耀光照…」（啟 21：23），唯一的道路就是主耶穌基督。
\subsubsection{結/約1:18/2:19－21/13:31/14:9/17:1－6}
\begin{table}[H]
\centering
\begin{tabular}{p{2.cm}|p{15.50cm}}\hline\hline
\multicolumn{1}{c|}{经文}&\multicolumn{1}{c}{内容}\\\hline
出34:6－7&耶和華在他面前宣告說，耶和華，耶和華，是有憐憫，有恩典的神，不輕易發怒，並有豐盛的慈愛和誠實。  為千萬人存留慈愛，赦免罪孽，過犯，和罪惡。萬不以有罪的為無罪，必追討他的罪…」\\\hline
出3:18－19&我要顯我一切的恩慈，在你面前經過，宣告我的名…」\\\hline

結&以西結在蒙召時得以見到神的面，看到祂的榮耀\\\hline
結&猶大國因長期悖逆叛道，在聖城被毀前，神的榮耀從城中離開，停在城東的山上\\\hline
結&神又給先知看到更新的聖城，神的榮耀再次在萬民中彰顯，「人就知道我是耶和華。」\\\hline
約1:18&從來沒有人看見神，只有在父懷裡的獨生子將祂表明出來\\\hline
約2:19－21&耶穌回答說：「你們拆毀這殿，我三日內要再建立起來。」 20 猶太人便說：「這殿是四十六年才造成的，你三日內就再建立起來嗎？」 21 但耶穌這話是以他的身體為殿。\\\hline
約13:31&如今人子得了榮耀，神在人子身上也得了榮耀\\\hline
約14:9&看見了我，就是看見了父\\\hline
約17:1－6&父阿，時候到了。願你榮耀你的兒子，使兒子也榮耀你。  正如你曾賜給他權柄，管理凡有血氣的，叫他將永生賜給你所賜給他的人。  認識你獨一的真神，並且認識你所差來的耶穌基督，這就是永生。我在地上已經榮耀你，你所託付我的事，我已成全了。父阿，現在求你使我同你享榮耀，就是未有世界以先，我同你所有的榮耀。  你從世上賜給我的人，我已將你的名顯明與他們。\\\hline
羅3:23&世人都犯了罪，虧缺了神的榮耀\\\hline
彼前2:12&你們在外邦人中，應當品行端正，叫那些毀謗你們是作惡的，因看見你們的好行為，便在鑒察的日子，歸榮耀給神。\\\hline
猶1:24&叫你們無瑕無疵，歡歡喜喜站在他榮耀之前…\\\hline
啟21:23&那城內又不用日月光照，因有神的榮耀光照…\\\hline

\end{tabular}
\end{table}



































